% Target: rmp-ii-tangent-projector
% Formula: P_T(A) = sum_i P_L^(i-1) tensor I_i tensor P_R^(i+1) - sum_i P_L^i tensor
%   P_R^(i+1).
% Ink: Canonical public tenkz and tenkzfree environments; no raw TikZ.
% Boundary: Each block closes only at its internal cup; the identity wire and all
%   displayed outer and physical endpoints stay open.
% Source: paper TN-Review-main.tex line 667; author source ImagesReview Section
%   II/ImagesReview Section II.tex lines 208-236
% Capabilities: grid, free-graph, typed-ports
\[
\begin{aligned}
  P_{T(A)} \;=\;& \sum_{i=1}^{n}\;
  % P_L^{(i-1)}: A_L window, all legs open, bracket (cup) at the east edge
  \begin{tenkz}[rows={ket:nopair, bra}, east=cup]
    \tndots & \tn[tri=l]{} & \tn[tri=l]{}\\
    \tndots & \tn*[tri=l]{} & \tn*[tri=l]{}
  \end{tenkz}
  \otimes
  \begin{tenkzfree}
    \tnput[boundary,ports={south:physical}]{it}{(0mm,5mm)}{}
    \tnput[boundary,ports={north:physical}]{ib}{(0mm,-5mm)}{}
    \tnjoin{it.south}{ib.north}
  \end{tenkzfree}
  \otimes
  % P_R^{(i+1)}: A_R window, bracket (cup) at the west edge
  \begin{tenkz}[rows={ket:nopair, bra}, west=cup]
    \tn[tri=r]{} & \tn[tri=r]{} & \tndots\\
    \tn*[tri=r]{} & \tn*[tri=r]{} & \tndots
  \end{tenkz}
  \\[0.4em]
  -\;& \sum_{i=1}^{n}\;
  % P_L^{(i)} (x) P_R^{(i+1)}: the two brackets meet back to back
  \begin{tenkz}[rows={ket:nopair, bra}, east=cup]
    \tndots & \tn[tri=l]{} & \tn[tri=l]{}\\
    \tndots & \tn*[tri=l]{} & \tn*[tri=l]{}
  \end{tenkz}
  \otimes
  \begin{tenkz}[rows={ket:nopair, bra}, west=cup]
    \tn[tri=r]{} & \tn[tri=r]{} & \tndots\\
    \tn*[tri=r]{} & \tn*[tri=r]{} & \tndots
  \end{tenkz}
\end{aligned}
\]
